% Einheit 5 – Anwendungen (Nr. 45–53)
\setcounter{aufgabe}{44}
\einheitenkopf[e5][5 Anwendungen]{Einheit 5 von 5 · Anwendungen}
\zweigzeile{Hier lernst du, Sachsituationen mit linearen Funktionen zu beschreiben und Tarife zu berechnen und zu vergleichen · neu in diesem Jahr · P10 oft · baut auf: Gleichung bestimmen und Schnittpunkt (Einheit 4), Funktionswerte (Einheit 3), Dreisatz}

\begin{aufgabe}{Ich erkenne im Text, was am Anfang da ist und was je Einheit dazukommt.\newline Unterstreiche den Anfangswert und kreise die Änderung je Einheit ein.}
\begin{teile}
\teil Ein Taxi kostet 4\,€ Grundpreis und 2,20\,€ je Kilometer.
\teil In einem Fass sind 80\,l Wasser. Je Minute laufen 5\,l ab.
\teil Ein Handyvertrag kostet 9,99\,€ je Monat und einmalig 25\,€ Anschlussgebühr.
\teil Lisa spart jede Woche 3\,€. Sie hat schon 45\,€.
\teil Eine Pflanze ist 12\,cm hoch und wächst jeden Tag 0,5\,cm.
\end{teile}
\end{aufgabe}

\verfahren{Gleichung zu einem Sachverhalt}

\begin{aufgabe}{Ich kann einem Tarif die passende Gleichung zuordnen.\newline Kreuze die passende Gleichung an. $y$ ist immer der Preis in €.}
\begin{teile}
\teil Fahrradverleih: 6\,€ Grundgebühr und 2\,€ je Stunde; $x$: Stunden \\
\kreuz{$y = 6x + 2$} \quad \kreuz{$y = 2x + 6$} \quad \kreuz{$y = 8x$}
\teil Schwimmbad: Jahreskarte 15\,€ und 3\,€ je Besuch; $x$: Besuche \\
\kreuz{$y = 18x$} \quad \kreuz{$y = 15x + 3$} \quad \kreuz{$y = 3x + 15$}
\teil Strom: 12\,€ Grundpreis im Monat und 0,30\,€ je kWh; $x$: kWh \\
\kreuz{$y = 0{,}3x + 12$} \quad \kreuz{$y = 12x + 0{,}3$} \quad \kreuz{$y = 12{,}3x$}
\teil Parkhaus: 2\,€ bei der Einfahrt und 5\,Cent je Minute; $x$: Minuten (P10 2021 OS) \\
\kreuz{$y = 5x + 2$} \quad \kreuz{$y = 2x + 0{,}05$} \quad \kreuz{$y = 0{,}05x + 2$} \quad \kreuz{$y = 0{,}5x + 2$}
\end{teile}
\end{aufgabe}

\begin{aufgabe}{Ich kann zu einem Sachverhalt die Gleichung aufstellen und einen Wert berechnen.\newline Stelle die Gleichung auf und berechne.}
\begin{teile}
\teil Ein Kletterpark kostet 8\,€ Eintritt und 4\,€ je Stunde. $x$: Stunden, $y$: Preis in €. \\ $y = $ \leerfeld
\teil Was kostet ein Besuch von 3 Stunden? \quad \leerfeld[€]
\teil Auf einem Konto sind 250\,€. Jede Woche werden 15\,€ abgehoben. $x$: Wochen, $y$: Kontostand in €. \\ $y = $ \leerfeld
\teil Wie hoch ist der Kontostand nach 8 Wochen? \quad \leerfeld[€]
\teil Ein Fitnessstudio kostet 60\,€ Aufnahmegebühr und 35\,€ im Monat. Stelle die Gleichung auf. Nach wie vielen Monaten hat Lea insgesamt mehr als 500\,€ bezahlt? (P10 2022 OS)
\end{teile}
\end{aufgabe}

\begin{aufgabe}{Ich kann zu einer Gleichung eine passende Situation beschreiben.\newline Beschreibe eine passende Situation und sage, was die beiden Zahlen bedeuten.}
\begin{teile}
\teil $y = 2{,}5x + 10$; \ $x$: Stunden, $y$: Preis in €
\teil $y = -3x + 60$; \ $x$: Minuten, $y$: Wassermenge in Litern
\end{teile}
\end{aufgabe}

\verfahren{Tarife vergleichen}

\begin{aufgabe}{Ich kann einem Tarif den passenden Graphen zuordnen.\newline Eine Kletterhalle hat drei Tarife. Tarif Start: 20\,€ im Monat und 5\,€ je Besuch. Tarif Einzel: 10\,€ je Besuch. Tarif Flat: 50\,€ im Monat, beliebig viele Besuche. Ordne jedem Tarif seinen Graphen zu.}

\noindent\begin{minipage}[t]{0.42\linewidth}\vspace{0pt}
\begin{ksys}[xmin=0,xmax=8,ymin=0,ymax=90,ystep=10,ablesen,xlabel=$x$ Besuche,ylabel=$y$ in €]
\gerade[1]{0}{50}{\mathrm{I}}
\gerade[6.5]{10}{0}{\mathrm{II}}
\gerade[7]{5}{20}{\mathrm{III}}
\end{ksys}
\end{minipage}\hfill
\begin{minipage}[t]{0.55\linewidth}\vspace{0pt}
\begin{teile}
\teil Tarif Einzel: \quad Graph \leerfeld
\teil Tarif Flat: \quad Graph \leerfeld
\teil Tarif Start: \quad Graph \leerfeld \\ Begründe deine Wahl mit zwei Eigenschaften des Graphen. (P10 2023 OS)
\end{teile}
\end{minipage}
\end{aufgabe}

\begin{aufgabe}{Ich kann zwei Tarife vergleichen und mich entscheiden.\newline Carsharing: Anbieter Nord verlangt 9\,€ je Fahrt und 0,30\,€ je km. Anbieter Süd verlangt nur 0,45\,€ je km.}
\begin{teile}
\teil Was kostet eine Fahrt von 40\,km bei Nord und bei Süd? Welcher Anbieter ist günstiger?
\teil Was kostet eine Fahrt von 80\,km bei Nord und bei Süd? Welcher Anbieter ist günstiger?
\teil Ab welcher Strecke ist Nord günstiger? Stelle für beide Anbieter die Gleichung auf und berechne. (P10 2016 OS)
\end{teile}
\end{aufgabe}

\begin{aufgabe}{Ich kann Tarife mit Freiminuten berechnen und vergleichen.\newline Ein E-Roller hat zwei Tarife. Basis: 2\,€ je Fahrt, die ersten 10 Minuten frei, danach 0,25\,€ je Minute. Plus: 5\,€ je Fahrt, die ersten 30 Minuten frei, danach 0,15\,€ je Minute.}
\begin{teile}
\teil Was kostet eine Fahrt von 8 Minuten mit Basis? \quad \leerfeld[€]
\teil Was kostet eine Fahrt von 40 Minuten mit Basis? \quad \leerfeld[€]
\teil Welcher Tarif ist für eine Fahrt von 40 Minuten günstiger? Begründe mit einer Rechnung. (P10 2023 OS)
\end{teile}
\end{aufgabe}

\verfahren{Prüfen, begründen, anwenden}

\begin{aufgabe}{Ich finde den Fehler beim Aufstellen einer Tarifgleichung.\newline Finde den Fehler, benenne ihn und korrigiere die Gleichung.}
\begin{teile}
\teil Parkplatz: 1,50\,€ bei der Einfahrt und 20\,Cent je Minute; $x$: Minuten, $y$: Preis in €. \\ Emma schreibt: \ $y = 20x + 1{,}5$
\teil Musikschule: 30\,€ Anmeldung und 45\,€ je Monat; $x$: Monate, $y$: Preis in €. \\ Paul schreibt: \ $y = 30x + 45$
\end{teile}
\end{aufgabe}

\begin{aufgabe}{Ich kann begründen, ob ein Tarif immer günstiger ist.\newline Entscheide ohne Rechnung und begründe.}
\begin{teile}
\teil Bei zwei Tarifen kostet jeder Kilometer gleich viel. Der erste Tarif hat die kleinere Grundgebühr. Ist er immer günstiger?
\teil Der erste Tarif hat die kleinere Grundgebühr, aber jeder Kilometer ist teurer als beim zweiten. Leon sagt: „Der erste ist immer günstiger.“ Stimmt das?
\end{teile}
\end{aufgabe}
